% Raport semestralny II — Dog FACS Dataset
% Kompilacja: tectonic Raport_DogFACS_2_semestr.tex
\documentclass[11pt,a4paper]{article}

\usepackage{fontspec}
\setmainfont{lmroman10-regular.otf}[
  BoldFont = lmroman10-bold.otf,
  ItalicFont = lmroman10-italic.otf,
  BoldItalicFont = lmroman10-bolditalic.otf,
]
\usepackage{polyglossia}
\setmainlanguage{polish}
\usepackage[margin=2.5cm]{geometry}
\usepackage{booktabs}
\usepackage{array}
\usepackage{xcolor}
\usepackage{titlesec}
\usepackage{fancyhdr}
\usepackage{lastpage}
\usepackage{enumitem}
\usepackage{hyperref}

\definecolor{sectionblue}{RGB}{0,90,160}
\definecolor{okgreen}{RGB}{0,140,60}

\titleformat{\section}{\Large\bfseries\color{sectionblue}}{\thesection.}{0.6em}{}
\titleformat{\subsection}{\large\bfseries}{\thesubsection.}{0.6em}{}

\pagestyle{fancy}
\fancyhf{}
\fancyhead[L]{Dog FACS Dataset}
\fancyhead[R]{Raport semestralny II}
\fancyfoot[C]{\thepage\ / \pageref{LastPage}}
\renewcommand{\headrulewidth}{0.4pt}
\renewcommand{\footrulewidth}{0.4pt}

\newcommand{\done}{\textcolor{okgreen}{Zrealizowane}}
\newcommand{\uzup}[1]{\textbf{[#1]}}

\begin{document}

% ============================ STRONA TYTUŁOWA ============================
\begin{titlepage}
\centering
\vspace*{2.5cm}

{\Large\textbf{Politechnika Gdańska}}\\[0.35em]
{\large Wydział Elektroniki, Telekomunikacji i Informatyki}\\[0.3em]
{\large Katedra Systemów Decyzyjnych i Robotyki}

\vspace{2.2cm}

{\Huge\textbf{Raport Semestralny}}\\[0.5em]
{\Large Semestr II (letni 2025/2026)}

\vspace{2.2cm}

{\LARGE\textbf{Dog FACS Dataset}}\\[0.5em]
{\large Utworzenie i Adnotacja Zbioru Danych\\do Analizy Emocji Psów}

\vspace{3cm}

\begin{tabular}{r l}
\textbf{Zespół projektowy:} & Danylo Lohachov (koordynator)\\
                            & Anton Shkrebela\\
                            & Danylo Zherzdiev\\
                            & Mariia Volkova\\[0.8em]
\textbf{Opiekun projektu:}  & dr hab. inż. Michał Czubenko\\
\end{tabular}

\vfill
Gdańsk, czerwiec 2026
\end{titlepage}

% ============================ SPIS TREŚCI ============================
\tableofcontents
\newpage

% ============================ 1. ZAŁOŻENIA ============================
\section{Założenia projektu}

\subsection{Cel główny}
Celem projektu jest stworzenie wysokiej jakości, publicznie dostępnego zbioru danych do
analizy emocji psów w formacie COCO, wspieranego przez dedykowany pipeline AI zdolny
do automatycznej adnotacji.

\subsection{Założenia na semestr II}
Zgodnie z planem przedstawionym w raporcie za semestr I, na semestr II zaplanowano:
\begin{enumerate}[leftmargin=3.2em]
  \item \textbf{Zbieranie danych} -- pobranie filmów psów z YouTube
  \item \textbf{Automatyczna adnotacja} -- przetworzenie klatek wideo przez pipeline AI
  \item \textbf{Weryfikacja manualna} -- ręczna weryfikacja próbki klatek
  \item \textbf{Klasyfikator emocji oparty na sieci neuronowej} -- zastąpienie rule-based
        classifiera modelem ML
  \item \textbf{Rozwój aplikacji do ręcznej adnotacji} -- pełny interfejs do manualnej
        adnotacji (bounding boxes, keypoints, emocje) z weryfikacją i korektą
  \item \textbf{Finalizacja datasetu} -- konwersja do formatu COCO, statystyki, dokumentacja
\end{enumerate}

% ============================ 2. REALIZACJA ============================
\section{Zrealizowany zakres prac}

\subsection{Podsumowanie realizacji}
Poniżej przedstawiono szczegółowy status każdego zadania semestru II:

\begin{table}[h!]
\centering
\begin{tabular}{p{8.2cm} l}
\toprule
\textbf{Zadanie} & \textbf{Status} \\
\midrule
Aplikacja webowa do ręcznej adnotacji (FastAPI + React) & \done \\
Ulepszenie modelu keypoints (detekcja dwuprzebiegowa) & \done \\
Korekta schematu keypoints (kanoniczna kolejność DogFLW) & \done \\
Naprawa klasyfikacji ras (mapowanie klas modelu) & \done \\
Stabilizacja Action Units (odporność na szum) & \done \\
Skrypty pobierania wideo (YouTube) i anotacji wsadowej & \done \\
Polonizacja interfejsu aplikacji & \done \\
Klasyfikator emocji (sieć neuronowa) & zmiana podejścia -- patrz sekcja 3.3 \\
\bottomrule
\end{tabular}
\caption{Status realizacji zadań semestru II}
\end{table}

\subsection{Produkty końcowe semestru II}
\begin{itemize}
  \item \textbf{Aplikacja webowa} -- pełnoprawne narzędzie do anotacji i weryfikacji
        (backend FastAPI + SessionStore, frontend React + TypeScript):
  \begin{itemize}
    \item przetwarzanie wgranego wideo przez pipeline (klatka neutralna, klatki
          szczytowe, Action Units, emocje),
    \item edytor 46 punktów kluczowych (przeciąganie, przełączanie widoczności),
    \item panel 21 Action Units z ręczną korektą intensywności,
    \item korekta rasy i emocji, eksport do COCO JSON,
    \item interfejs w całości w języku polskim.
  \end{itemize}
  \item \textbf{Ulepszony pipeline AI} -- istotne poprawki jakości anotacji
        (szczegóły w sekcji 3).
  \item \textbf{Narzędzia danych} -- skrypty pobierania wideo z YouTube (yt-dlp),
        anotacji wsadowej oraz skrypty diagnostyczne pipeline'u.
  \item \textbf{Repozytorium GitHub} -- kod źródłowy z CI/CD.
  \item \textbf{Dokumentacja i prezentacja} -- dokumentacja sprintów, prezentacja
        końcowa projektu (PDF/PPTX).
\end{itemize}

% ============================ 3. ZMIANY ============================
\section{Zmiany względem pierwotnego planu}

W trakcie semestru II wprowadzono trzy istotne zmiany:

\subsection{Zmiana 1: Korekta schematu punktów kluczowych (DogFLW)}
\begin{itemize}
  \item \textbf{Problem:} wykryto, że przyjęty w semestrze I schemat nazw i kolejności
        46 punktów kluczowych nie odpowiadał rzeczywistej kolejności kanałów wyjściowych
        modelu -- punkty były semantycznie błędnie przypisane (np. punkty
        ,,oczu'' lądowały na policzku), a Action Units liczone były na złych punktach.
  \item \textbf{Rozwiązanie:} schemat przepisano na kanoniczną kolejność
        \textbf{DogFLW} (Dog Facial Landmarks in the Wild, arXiv:2405.11501),
        co potwierdzono empirycznie (uśredniony szablon pozycji kanałów na 66 obrazach).
        Dodatkowo wprowadzono \textbf{dwuprzebiegową detekcję} (zgrubna na całym psie,
        dokładna na przyciętym pysku).
  \item \textbf{Wynik:} punkty trafiają na właściwe cechy anatomiczne, a AU liczone są
        na poprawnych punktach.
\end{itemize}

\subsection{Zmiana 2: Naprawa mapowania klas rasy}
\begin{itemize}
  \item \textbf{Problem:} model klasyfikacji ras zwracał z wysoką pewnością błędne rasy --
        plik etykiet miał inną kolejność klas niż ta, w której model był trenowany.
  \item \textbf{Rozwiązanie:} prawdziwe mapowanie indeks $\rightarrow$ rasa odtworzono
        empirycznie (test modelu na próbkach wielu ras) i przebudowano plik etykiet.
  \item \textbf{Wynik:} poprawa z 0/6 do \textbf{6/6} poprawnych ras na próbie testowej.
\end{itemize}

\subsection{Zmiana 3: Klasyfikator emocji -- stabilizacja zamiast sieci neuronowej}
\begin{itemize}
  \item \textbf{Plan:} zastąpienie klasyfikatora rule-based siecią neuronową.
  \item \textbf{Realizacja:} zachowano klasyfikator regułowy DogFACS, koncentrując się
        na \textbf{stabilizacji obliczeń Action Units}: przycięcie wartości stosunków
        (odporność na degenerację klatki neutralnej) oraz bramkowanie aktywacji AU
        pewnością punktów kluczowych (eliminacja fałszywych aktywacji z szumu).
  \item \textbf{Powód:} po korekcie schematu DogFLW konieczne było najpierw zapewnienie
        wiarygodnych wartości AU -- trening sieci na błędnie liczonych AU utrwaliłby
        wady danych. Stabilny, interpretowalny klasyfikator regułowy stanowi solidną
        bazę do przyszłego treningu sieci na zweryfikowanych danych.
\end{itemize}

% ============================ 4. WKŁAD ============================
\section{Wkład poszczególnych członków zespołu}

\subsection{Danylo Lohachov (U1) -- Koordynator projektu}
\begin{itemize}
  \item Koordynacja prac zespołu i harmonogramu
  \item Rozwój frontendu aplikacji webowej (edytor keypoints, panel AU, eksport)
  \item Polonizacja interfejsu aplikacji
  \item Dokumentacja projektowa i prezentacja końcowa
  \item Zarządzanie repozytorium GitHub i Linear
\end{itemize}
\textbf{Główne osiągnięcia:}
\begin{itemize}
  \item Działająca aplikacja webowa do anotacji (frontend)
  \item Kompletna dokumentacja i prezentacja końcowa projektu
\end{itemize}

\subsection{Anton Shkrebela (U2) -- Specjalista AI/ML}
\begin{itemize}
  \item Diagnoza i korekta schematu punktów kluczowych (DogFLW)
  \item Implementacja dwuprzebiegowej detekcji keypoints
  \item Stabilizacja obliczeń Action Units (clamp, bramkowanie pewnością)
  \item Ulepszenie selektora klatek szczytowych (adaptacyjna separacja)
\end{itemize}
\textbf{Główne osiągnięcia:}
\begin{itemize}
  \item Punkty kluczowe trafiające na właściwe cechy anatomiczne pyska
  \item Wiarygodne wartości AU (eliminacja fałszywych aktywacji)
\end{itemize}

\subsection{Danylo Zherzdiev (U3) -- Backend Developer}
\begin{itemize}
  \item Backend aplikacji webowej (FastAPI, SessionStore, REST API)
  \item Naprawa mapowania klas modelu ras (0/6 $\rightarrow$ 6/6)
  \item Integracja ulepszeń w zunifikowanym pipeline
  \item Eksport COCO z rozszerzeniami DogFACS
\end{itemize}
\textbf{Główne osiągnięcia:}
\begin{itemize}
  \item Stabilny backend z obsługą sesji anotacyjnych
  \item Poprawna klasyfikacja ras (6/6 na próbie testowej)
\end{itemize}

\subsection{Mariia Volkova (U4) -- Data Engineer}
\begin{itemize}
  \item Skrypty pobierania wideo z YouTube (yt-dlp) i ich naprawa
  \item Anotacja wsadowa klatek przez pipeline
  \item Ręczna weryfikacja anotacji w aplikacji webowej
  \item Przygotowanie zapytań wyszukiwania dla zbierania danych
\end{itemize}
\textbf{Główne osiągnięcia:}
\begin{itemize}
  \item Działający tor pozyskiwania danych (wyszukiwanie--pobranie--anotacja)
\end{itemize}

\subsection{Podsumowanie wkładu}

\begin{table}[h!]
\centering
\begin{tabular}{l p{7.2cm} c}
\toprule
\textbf{Osoba} & \textbf{Główny obszar odpowiedzialności} & \textbf{Wkład} \\
\midrule
Danylo Lohachov  & Koordynacja, frontend, dokumentacja      & 30\% \\
Anton Shkrebela  & Keypoints, Action Units, DogFACS         & 30\% \\
Danylo Zherzdiev & Backend, rasy, pipeline, COCO            & 20\% \\
Mariia Volkova   & Dane, anotacja wsadowa, weryfikacja      & 20\% \\
\bottomrule
\end{tabular}
\caption{Podział wkładu w projekcie (semestr II)}
\end{table}

% ============================ 5. OCENY ============================
\section{Proponowane oceny}

Na podstawie zrealizowanego zakresu prac, osiągniętych wyników i zaangażowania
poszczególnych członków zespołu, proponujemy następujące oceny:

\begin{table}[h!]
\centering
\begin{tabular}{l c p{8cm}}
\toprule
\textbf{Osoba} & \textbf{Ocena} & \textbf{Uzasadnienie} \\
\midrule
Danylo Lohachov  & 5.0 & Skuteczna koordynacja, frontend aplikacji anotacyjnej,
                         kompletna dokumentacja i prezentacja końcowa \\[0.5em]
Anton Shkrebela  & 5.0 & Diagnoza i naprawa krytycznego błędu schematu keypoints,
                         dwuprzebiegowa detekcja, stabilizacja AU \\[0.5em]
Danylo Zherzdiev & 4.0 & Backend aplikacji, naprawa klasyfikacji ras,
                         integracja pipeline'u \\[0.5em]
Mariia Volkova   & 4.0 & Tor pozyskiwania danych, anotacja wsadowa,
                         weryfikacja manualna \\
\bottomrule
\end{tabular}
\caption{Proponowane oceny za semestr II}
\end{table}

\subsection{Uzasadnienie ocen}
\begin{enumerate}[leftmargin=3.2em]
  \item \textbf{Jakość anotacji} -- wykryto i naprawiono krytyczne błędy
        (schemat keypoints, mapowanie ras), które bezpośrednio decydują
        o wartości naukowej zbioru danych
  \item \textbf{Kompletny produkt} -- działająca aplikacja webowa obejmująca pełen
        cykl: przetwarzanie wideo, przegląd, korektę i eksport COCO
  \item \textbf{Rzetelność inżynierska} -- zmiany potwierdzane empirycznie
        (testy na próbach, porównania przed/po)
  \item \textbf{Współpraca zespołowa} -- spójny podział ról kontynuowany
        z semestru I, wspólne rozwiązywanie problemów
  \item \textbf{Adaptacja do zmian} -- świadoma decyzja o stabilizacji AU przed
        treningiem sieci neuronowej, chroniąca jakość przyszłego datasetu
\end{enumerate}

% ============================ 6. PODSUMOWANIE ============================
\section{Podsumowanie projektu i możliwości dalszego rozwoju}

Projekt zakończył się dostarczeniem działającego systemu \textbf{end-to-end}:
pipeline AI (detekcja psa, klasyfikacja rasy, 46 punktów kluczowych DogFLW,
21 Action Units, klasyfikacja emocji) wraz z aplikacją webową do ręcznej anotacji
i weryfikacji oraz eksportem do formatu COCO.

Możliwe kierunki dalszego rozwoju:
\begin{enumerate}[leftmargin=3.2em]
  \item \textbf{Dokończenie zbioru danych} -- kontynuacja zbierania i weryfikacji danych
  \item \textbf{Sieć neuronowa AU} -- trening MLP (138 wejść $\rightarrow$ 21 AU)
        na zweryfikowanych danych, w miejsce klasyfikatora regułowego
  \item \textbf{Fine-tuning modelu keypoints} -- doszkolenie na trudnych przypadkach
        (profile, rasy o płaskich pyskach)
  \item \textbf{Publikacja datasetu} -- udostępnienie zbioru wraz ze statystykami
        i dokumentacją
\end{enumerate}

\vspace{1.5em}
\noindent\rule{\textwidth}{0.4pt}
\begin{center}
\textbf{Dokument przygotowany przez zespół projektowy Dog FACS Dataset}\\
Gdańsk, czerwiec 2026
\end{center}

\end{document}
